\documentclass[12pt,a4paper]{ctexart}

\usepackage{amsmath,amssymb,amsfonts}
\usepackage{geometry}
\usepackage{booktabs}
\usepackage{enumitem}
\usepackage{setspace}
\usepackage{graphicx}
\usepackage{booktabs}

\geometry{left=2.5cm,right=2.5cm,top=2.5cm,bottom=2.5cm}
\setstretch{1.3}

\title{问题一第（1）种情形的优化模型建立与线性化求解}
\author{}
\date{}

\begin{document}

\maketitle

\section{问题分析}

问题一第（1）种情形要求在 2024--2030 年期间，为乡村现有耕地和大棚制定最优种植方案。
题目假定各种农作物未来的预期销售量、种植成本、亩产量和销售价格相对于 2023 年保持稳定，
并且每季种植的农作物在当季销售。若某种作物每季的总产量超过相应的预期销售量，
则超过部分滞销，造成浪费，不产生销售收入。

因此，本问题的核心是在满足地块类型适宜性、种植面积、轮作、豆类作物种植要求以及田间管理要求的前提下，
最大化 2024--2030 年的总收益。由于目标函数中涉及实际销售量
$\min\{Q_{tsi},D_{si}\}$，因此原始模型可建立为混合整数非线性规划模型。

\section{集合与下标}

为便于建模，定义如下集合与下标：

\begin{itemize}[leftmargin=2em]
    \item $T=\{2024,2025,\ldots,2030\}$：规划年份集合；
    \item $L$：全部种植单元集合，包括露天地块和大棚；
    \item $L_F$：平旱地集合；
    \item $L_T$：梯田集合；
    \item $L_S$：山坡地集合；
    \item $L_I$：水浇地集合；
    \item $L_N$：普通大棚集合；
    \item $L_H$：智慧大棚集合；
    \item $C=\{1,2,\ldots,41\}$：全部农作物集合；
    \item $C_g$：粮食类作物集合，不包括水稻；
    \item $C_r$：水稻作物集合；
    \item $C_{v1}$：第一季可种植蔬菜集合；
    \item $C_{v2}$：水浇地第二季可种植蔬菜集合，即大白菜、白萝卜、红萝卜；
    \item $C_m$：食用菌集合；
    \item $C_b$：豆类作物集合；
    \item $S=\{0,1,2\}$：种植季次集合，其中 $s=0$ 表示单季，$s=1$ 表示第一季，$s=2$ 表示第二季。
\end{itemize}

其中，根据题目给出的作物编号，可取：
\[
C_g=\{1,2,\ldots,15\},\quad
C_r=\{16\},
\]
\[
C_{v1}=\{17,18,\ldots,34\},\quad
C_{v2}=\{35,36,37\},
\]
\[
C_m=\{38,39,40,41\},
\]
\[
C_b=\{1,2,3,4,5,17,18,19\}.
\]

\section{参数定义}

定义如下模型参数：

\begin{itemize}[leftmargin=2em]
    \item $A_l$：地块或大棚 $l$ 的面积，单位为亩；
    \item $Y_{lsi}$：作物 $i$ 在地块 $l$、季次 $s$ 的亩产量，单位为斤/亩；
    \item $K_{lsi}$：作物 $i$ 在地块 $l$、季次 $s$ 的种植成本，单位为元/亩；
    \item $P_{si}$：作物 $i$ 在季次 $s$ 的销售单价，单位为元/斤；
    \item $D_{si}$：作物 $i$ 在季次 $s$ 的预期销售量，单位为斤；
    \item $E_{lsi}$：作物适宜性参数，若作物 $i$ 可在地块 $l$、季次 $s$ 种植，则 $E_{lsi}=1$，否则 $E_{lsi}=0$；
    \item $M_i$：作物 $i$ 在同一年同一季最多允许分布的地块数量；
    \item $m_{li}$：作物 $i$ 在地块 $l$ 上的最小种植面积；
    \item $x^{2023}_{lsi}$：2023 年作物 $i$ 在地块 $l$、季次 $s$ 的实际种植面积；
    \item $z^{2023}_{lsi}$：2023 年作物 $i$ 是否在地块 $l$、季次 $s$ 种植；
    \item $B^{2023}_l$：2023 年地块 $l$ 上豆类作物的实际种植面积。
\end{itemize}

若附件中销售价格为区间 $[P^{\min}_{si},P^{\max}_{si}]$，则可取区间均值作为销售价格：
\[
P_{si}=\frac{P^{\min}_{si}+P^{\max}_{si}}{2}.
\]

由于问题一假定未来各项参数相对于 2023 年保持稳定，因此可用 2023 年该作物的总产量近似表示预期销售量：
\[
D_{si}=\sum_{l\in L}Y_{lsi}x^{2023}_{lsi}.
\]

\section{决策变量}

定义如下决策变量：

\[
x_{tlsi}\geq 0,
\]

表示第 $t$ 年、地块 $l$、季次 $s$ 种植作物 $i$ 的面积，单位为亩。

\[
z_{tlsi}\in\{0,1\},
\]

表示第 $t$ 年、地块 $l$、季次 $s$ 是否种植作物 $i$。若种植，则 $z_{tlsi}=1$；否则 $z_{tlsi}=0$。

对于水浇地，定义模式选择变量：
\[
r_{tl}\in\{0,1\},\quad l\in L_I.
\]

其中，$r_{tl}=1$ 表示水浇地 $l$ 在第 $t$ 年选择单季种植水稻；
$r_{tl}=0$ 表示水浇地 $l$ 在第 $t$ 年选择两季种植蔬菜。

定义总产量变量：
\[
Q_{tsi}\geq 0,
\]

表示第 $t$ 年、第 $s$ 季、作物 $i$ 的总产量。

\section{目标函数}

在第（1）种情形下，若某种作物的总产量超过预期销售量，则超过部分滞销，不能产生正常销售收入。
因此，作物 $i$ 在第 $t$ 年、第 $s$ 季的实际可销售产量为：
\[
\min\{Q_{tsi},D_{si}\}.
\]

该作物的销售收入为：
\[
P_{si}\min\{Q_{tsi},D_{si}\}.
\]

总种植成本为：
\[
\sum_{t\in T}\sum_{l\in L}\sum_{s\in S}\sum_{i\in C}
K_{lsi}x_{tlsi}.
\]

因此，以 2024--2030 年总利润最大为目标，可建立目标函数：
\[
\max Z=
\sum_{t\in T}\sum_{s\in S}\sum_{i\in C}
P_{si}\min\{Q_{tsi},D_{si}\}
-
\sum_{t\in T}\sum_{l\in L}\sum_{s\in S}\sum_{i\in C}
K_{lsi}x_{tlsi}.
\]

其中，总产量满足：
\[
Q_{tsi}=\sum_{l\in L}Y_{lsi}x_{tlsi},
\quad \forall t\in T,\ s\in S,\ i\in C.
\]

\section{约束条件}

\subsection{作物适宜性约束}

若作物 $i$ 不适宜在地块 $l$、季次 $s$ 种植，则其种植面积应为 0。
用适宜性参数 $E_{lsi}$ 表示为：
\[
x_{tlsi}\leq A_lE_{lsi}z_{tlsi},
\quad \forall t\in T,\ l\in L,\ s\in S,\ i\in C.
\]

该约束同时保证：当 $E_{lsi}=0$ 时，$x_{tlsi}=0$；当 $z_{tlsi}=0$ 时，$x_{tlsi}=0$。

\subsection{最小种植面积约束}

为避免某种作物在单个地块上的种植面积过小，引入最小种植面积约束：
\[
x_{tlsi}\geq m_{li}z_{tlsi},
\quad \forall t\in T,\ l\in L,\ s\in S,\ i\in C.
\]

其中，$m_{li}$ 可根据地块面积和实际管理需要确定。例如可取：
\[
m_{li}=\min\{0.3,0.2A_l\}.
\]

\subsection{平旱地、梯田和山坡地约束}

平旱地、梯田和山坡地每年只能种植一季粮食类作物，因此有：
\[
\sum_{i\in C_g}x_{tl0i}=A_l,
\quad \forall t\in T,\ l\in L_F\cup L_T\cup L_S.
\]

同时，这三类地块不安排第一季和第二季作物：
\[
x_{tl1i}=0,\quad x_{tl2i}=0,
\quad \forall t\in T,\ l\in L_F\cup L_T\cup L_S,\ i\in C.
\]

\subsection{水浇地种植模式约束}

水浇地有两种种植模式：一是单季种植水稻，二是两季种植蔬菜。

若选择单季种植水稻，则有：
\[
\sum_{i\in C_r}x_{tl0i}=A_lr_{tl},
\quad \forall t\in T,\ l\in L_I.
\]

若选择两季种植蔬菜，则第一季满足：
\[
\sum_{i\in C_{v1}}x_{tl1i}=A_l(1-r_{tl}),
\quad \forall t\in T,\ l\in L_I.
\]

第二季满足：
\[
\sum_{i\in C_{v2}}x_{tl2i}=A_l(1-r_{tl}),
\quad \forall t\in T,\ l\in L_I.
\]

同时，水浇地第二季只能在大白菜、白萝卜和红萝卜中选择一种作物：
\[
\sum_{i\in C_{v2}}z_{tl2i}=1-r_{tl},
\quad \forall t\in T,\ l\in L_I.
\]

\subsection{普通大棚约束}

普通大棚每年第一季种植蔬菜，第二季种植食用菌，因此：
\[
\sum_{i\in C_{v1}}x_{tl1i}=A_l,
\quad \forall t\in T,\ l\in L_N.
\]

\[
\sum_{i\in C_m}x_{tl2i}=A_l,
\quad \forall t\in T,\ l\in L_N.
\]

普通大棚不安排单季作物：
\[
x_{tl0i}=0,
\quad \forall t\in T,\ l\in L_N,\ i\in C.
\]

\subsection{智慧大棚约束}

智慧大棚每年可种植两季蔬菜，因此：
\[
\sum_{i\in C_{v1}}x_{tl1i}=A_l,
\quad \forall t\in T,\ l\in L_H.
\]

\[
\sum_{i\in C_{v1}}x_{tl2i}=A_l,
\quad \forall t\in T,\ l\in L_H.
\]

智慧大棚不安排单季作物：
\[
x_{tl0i}=0,
\quad \forall t\in T,\ l\in L_H,\ i\in C.
\]

\subsection{不连续重茬约束}

根据题意，每种作物在同一地块不能连续重茬种植。

对于一年一季的平旱地、梯田和山坡地，有：
\[
z_{tl0i}+z_{t-1,l,0,i}\leq 1,
\quad \forall t\in T,\ l\in L_F\cup L_T\cup L_S,\ i\in C.
\]

当 $t=2024$ 时，$z_{2023,l,0,i}$ 由附件 2 中 2023 年的实际种植情况给出。

对于一年两季的水浇地和大棚，同一年相邻两季之间不能重茬：
\[
z_{tl1i}+z_{tl2i}\leq 1,
\quad \forall t\in T,\ l\in L_I\cup L_N\cup L_H,\ i\in C.
\]

相邻年份之间，上一年第二季与下一年第一季也不能重茬：
\[
z_{tl2i}+z_{t+1,l,1,i}\leq 1,
\quad \forall t=2024,\ldots,2029,\ l\in L_I\cup L_N\cup L_H,\ i\in C.
\]

对于 2024 年第一季，还需与 2023 年最后一季的种植情况衔接：
\[
z_{2024,l,1,i}+z^{2023}_{l,\mathrm{last},i}\leq 1,
\quad \forall l\in L_I\cup L_N\cup L_H,\ i\in C.
\]

其中，$z^{2023}_{l,\mathrm{last},i}$ 表示 2023 年地块 $l$ 最后一季是否种植作物 $i$。

\subsection{三年内至少种植一次豆类作物约束}

题目要求从 2023 年开始，每个地块的所有土地三年内至少种植一次豆类作物。
因此，对于每个地块 $l$，每个三年滚动窗口，豆类作物种植面积累计应不少于该地块面积。

对于包含 2023 年的窗口，有：
\[
B^{2023}_l+
\sum_{t=2024}^{2025}
\sum_{s\in S}
\sum_{i\in C_b}
x_{tlsi}
\geq A_l,
\quad \forall l\in L.
\]

对于完全位于规划期内的三年窗口，有：
\[
\sum_{t=\tau}^{\tau+2}
\sum_{s\in S}
\sum_{i\in C_b}
x_{tlsi}
\geq A_l,
\quad \forall l\in L,\ \tau=2024,2025,\ldots,2028.
\]

该约束表示每块地的全部面积在任意连续三年内至少经历一次豆类作物种植，
从而满足轮作和改善土壤肥力的要求。

\subsection{作物种植不宜过分分散约束}

为方便田间管理，限制每种作物在同一年同一季分布的地块数量：
\[
\sum_{l\in L}z_{tlsi}\leq M_i,
\quad \forall t\in T,\ s\in S,\ i\in C.
\]

其中，$M_i$ 可根据 2023 年该作物的实际分布地块数量确定。例如，若作物 $i$ 在 2023 年分布于
$M_i^{2023}$ 个地块，则可取：
\[
M_i=\left\lceil 1.2M_i^{2023}\right\rceil.
\]

\subsection{变量取值约束}

\[
x_{tlsi}\geq 0,
\quad \forall t\in T,\ l\in L,\ s\in S,\ i\in C.
\]

\[
Q_{tsi}\geq 0,
\quad \forall t\in T,\ s\in S,\ i\in C.
\]

\[
z_{tlsi}\in\{0,1\},
\quad \forall t\in T,\ l\in L,\ s\in S,\ i\in C.
\]

\[
r_{tl}\in\{0,1\},
\quad \forall t\in T,\ l\in L_I.
\]

\section{原始模型汇总}

综上，问题一第（1）种情形的原始优化模型为：

\[
\max Z=
\sum_{t\in T}\sum_{s\in S}\sum_{i\in C}
P_{si}\min\{Q_{tsi},D_{si}\}
-
\sum_{t\in T}\sum_{l\in L}\sum_{s\in S}\sum_{i\in C}
K_{lsi}x_{tlsi},
\]

满足：

\[
Q_{tsi}=\sum_{l\in L}Y_{lsi}x_{tlsi},
\quad \forall t\in T,\ s\in S,\ i\in C,
\]

以及作物适宜性约束、面积约束、水浇地种植模式约束、大棚种植约束、
不连续重茬约束、三年内至少种植一次豆类作物约束、种植分散度约束和变量取值约束。

由于该模型目标函数中含有 $\min\{Q_{tsi},D_{si}\}$，因此属于混合整数非线性规划模型。
后续可通过引入实际销售量变量 $S_{tsi}$ 对其进行线性化处理，从而转化为混合整数线性规划模型，
便于使用 Gurobi 等商业优化软件求解。
\section{原始模型的智能优化算法求解}

\subsection{算法设计思路}

在原始优化模型中，决策变量包括连续种植面积变量 $x_{tlsi}$ 和是否种植变量 $z_{tlsi}$。
由于模型同时包含年份、地块、季次和作物等多个维度，变量数量较多，且目标函数中含有
$\min\{Q_{tsi},D_{si}\}$，直接使用传统精确算法求解存在一定复杂性。因此，本文首先采用智能优化算法对原始模型进行近似求解，并与后续 Gurobi 精确求解结果进行比较。

本文选取模拟退火算法和遗传算法两种智能优化算法。为了便于算法编码和提高搜索效率，
在智能优化算法中将一个完整的种植方案编码为：
\[
\Omega=\{m_{tl}\mid t\in T,\ l\in L\},
\]
其中 $m_{tl}$ 表示第 $t$ 年地块 $l$ 的种植模式。对于一年一季地块，
$m_{tl}$ 表示该地块当年种植的一种作物；对于水浇地、普通大棚和智慧大棚，
$m_{tl}$ 表示该地块当年第一季和第二季的作物组合。

需要说明的是，题目允许同一地块同一季合种不同作物。为了降低智能优化算法的编码复杂度，
本文在启发式算法中采用“每个地块每季选择一种主作物并种满该季面积”的方式进行搜索。
该处理相当于在原模型可行域的一个子空间内寻找较优方案，能够用于比较不同智能优化算法的搜索能力。
后续使用 Gurobi 求解时，再回到完整的连续面积变量模型。

\subsection{适应度函数}

对于给定的种植方案，先计算第 $t$ 年第 $s$ 季作物 $i$ 的总产量：
\[
Q_{tsi}=\sum_{l\in L}Y_{lsi}x_{tlsi}.
\]

在问题一第（1）种情形下，超过预期销售量的部分滞销，不能产生收入。因此实际销售收入为：
\[
R=\sum_{t\in T}\sum_{s\in S}\sum_{i\in C}
P_{si}\min\{Q_{tsi},D_{si}\}.
\]

总种植成本为：
\[
C=\sum_{t\in T}\sum_{l\in L}\sum_{s\in S}\sum_{i\in C}
K_{lsi}x_{tlsi}.
\]

原始利润为：
\[
\Pi=R-C.
\]

由于智能优化算法在搜索过程中可能产生违反约束的方案，本文引入罚函数，将违反约束的程度转化为惩罚项。设 $V_1$ 表示重茬约束违反次数，$V_2$ 表示三年内豆类作物种植约束违反次数，$V_3$ 表示作物分散度超限次数，则适应度函数设为：
\[
F=\Pi-\lambda_1V_1-\lambda_2V_2-\lambda_3V_3.
\]

其中，$\lambda_1,\lambda_2,\lambda_3$ 为惩罚系数。由于重茬约束和三年内豆类作物约束属于硬约束，
其惩罚系数取较大值；作物分散度主要体现田间管理便利性，本文将其作为软约束处理，惩罚系数相对较小。
在程序计算中，重茬约束和豆类约束的单次惩罚系数取 $10^8$，作物分散度超限的单次惩罚系数取 $2\times 10^5$。

\subsection{模拟退火算法}

模拟退火算法的基本思想是从一个初始种植方案出发，通过随机扰动不断产生新方案。
若新方案的适应度更高，则接受该方案；若新方案的适应度较低，则以一定概率接受该方案，
从而避免算法过早陷入局部最优。

本文模拟退火算法的主要步骤如下：

\begin{enumerate}
    \item 随机生成一个初始种植方案，并通过修复函数尽量满足重茬约束和三年内豆类作物约束；
    \item 计算当前方案的适应度函数值 $F$；
    \item 随机选择某一年某一地块，改变其种植模式，形成邻域解；
    \item 若邻域解适应度更高，则接受该解；若邻域解适应度更低，则以概率
    \[
    p=\exp\left(\frac{F_{\text{new}}-F_{\text{old}}}{T}\right)
    \]
    接受该解；
    \item 按照 $T\leftarrow \alpha T$ 更新温度，其中 $\alpha$ 为降温系数；
    \item 重复上述过程，直到达到最大迭代次数。
\end{enumerate}

在本次计算中，模拟退火算法最大迭代次数为 800 次，初始温度为 $10^7$，降温系数为 0.995。

\subsection{遗传算法}

遗传算法将一个完整的 2024--2030 年种植方案视为一个个体，将一组种植方案视为种群。
通过选择、交叉和变异操作，使种群不断向高适应度方向进化。

本文遗传算法的主要步骤如下：

\begin{enumerate}
    \item 随机生成初始种群，每个个体代表一个完整的 2024--2030 年种植方案；
    \item 计算每个个体的适应度函数值；
    \item 采用锦标赛选择方法，从种群中选择适应度较高的个体作为父代；
    \item 对两个父代方案进行交叉，随机继承不同年份和地块的种植模式，形成子代；
    \item 对子代进行随机变异，即随机改变某一年某一地块的作物种植模式；
    \item 对部分子代进行修复，使其尽量满足重茬约束和豆类轮作约束；
    \item 保留适应度最高的若干个体进入下一代，直到达到最大进化代数。
\end{enumerate}

在本次计算中，遗传算法种群规模为 30，最大进化代数为 20，变异概率为 0.15。

\subsection{计算结果}

两种智能优化算法的计算结果如表所示。

\begin{table}[htbp]
\centering
\caption{原始模型下两种智能优化算法的计算结果}
\label{tab:sa-ga-result}
\small
\setlength{\tabcolsep}{3pt}
\resizebox{\textwidth}{!}{%
\begin{tabular}{lrrrrrrc}
\toprule
算法 & 运行时间/s & 原始利润/元 & 总收入/元 & 总成本/元 & 滞销浪费/斤 & 惩罚项 & 是否有效 \\
\midrule
模拟退火 SA & 9.1436 & 9705157.00 & 16368205.00 & 6663048.00 & 9168905.00 & 9000000.00 & 是 \\
遗传算法 GA & 8.6847 & 10921219.75 & 17233973.75 & 6312754.00 & 7790400.00 & 7000000.00 & 是 \\
\bottomrule
\end{tabular}%
}
\end{table}

从表中可以看出，模拟退火算法和遗传算法均在较短时间内得到了满足主要硬约束的有效种植方案。
其中，模拟退火算法的运行时间为 9.1436 秒，原始利润为 9705157.00 元；
遗传算法的运行时间为 8.6847 秒，原始利润为 10921219.75 元。
遗传算法的原始利润高于模拟退火算法，且滞销浪费产量更少，说明遗传算法在本问题上的搜索效果更好。

从惩罚项看，两种算法均未违反重茬约束和三年内豆类作物约束，主要违反的是作物分散度约束。
模拟退火算法中作物分散度超限 45 次，遗传算法中作物分散度超限 35 次。
这说明两种算法虽然能够获得满足主要种植规律的方案，但在田间管理便利性方面仍有改进空间。
若进一步提高作物分散度惩罚系数，或在邻域搜索和遗传变异过程中增加集中种植修复操作，
可以进一步减少作物分散度超限情况。

\subsection{结果分析}

综合比较两种算法，遗传算法在本次实验中表现更优。其原因在于，遗传算法通过种群并行搜索，
能够同时保留多个较优方案，并通过交叉和变异组合不同地块和年份的种植安排，搜索空间覆盖范围较广；
而模拟退火算法主要依靠单个当前解的随机扰动进行搜索，虽然具有跳出局部最优的能力，
但在有限迭代次数下更容易受到初始解和邻域结构的影响。

因此，在原始模型的智能优化算法求解中，遗传算法得到的方案具有更高的利润、更低的惩罚项和更少的滞销浪费，
可作为后续与 Gurobi 精确求解结果进行比较的主要启发式算法结果。
\subsection{Gurobi 对原始模型的求解}

为进一步检验智能优化算法所得结果的有效性，本文使用商业优化软件 Gurobi 对原始模型进行求解。
与模拟退火算法和遗传算法不同，Gurobi 通过建立整数规划模型并调用分支定界、割平面等优化方法进行求解，
能够在给定最优性间隙范围内给出较为稳定的优化结果。

在原始模型中，问题一第（1）种情形的核心特征是超过预期销售量的部分滞销，不能产生销售收入。
因此，模型中仍保留实际销售量与总产量、预期销售量之间的最小值关系：
\[
S_{tsi}=\min\{Q_{tsi},D_{si}\}.
\]
其中，$Q_{tsi}$ 表示第 $t$ 年第 $s$ 季作物 $i$ 的总产量，$D_{si}$ 表示作物 $i$ 在第 $s$ 季的预期销售量，
$S_{tsi}$ 表示实际可销售产量。本文在 Gurobi 中使用一般约束表示上述最小值关系，从而保留原始模型的非线性结构。

受当前 Gurobi 许可证规模限制和计算规模影响，本文在 Gurobi 求解中采用与智能优化算法相同的种植模式编码方法。
即每个地块在每一年从若干个候选种植模式中选择一种，候选模式根据作物适宜性、2023 年实际种植模式、
需求修正后的收益水平以及豆类轮作要求生成。该处理能够在控制变量和约束数量的同时，保留作物适宜性、
不连续重茬、三年内至少种植一次豆类作物以及滞销浪费等核心条件。

设二元变量 $y_{tlk}$ 表示第 $t$ 年地块 $l$ 是否选择第 $k$ 个候选种植模式：
\[
y_{tlk}=
\begin{cases}
1, & \text{若第 }t\text{ 年地块 }l\text{ 选择第 }k\text{ 个候选模式},\\
0, & \text{否则}.
\end{cases}
\]

每个地块每年只能选择一个种植模式，因此有：
\[
\sum_{k\in M_l}y_{tlk}=1,
\quad \forall t\in T,\ l\in L.
\]

对于总产量，仍按作物、年份和季次汇总：
\[
Q_{tsi}=\sum_{l\in L}\sum_{k\in M_l}a_{lksi}Y_{lsi}A_l y_{tlk},
\]
其中，$a_{lksi}$ 表示地块 $l$ 的第 $k$ 个候选模式是否在季次 $s$ 种植作物 $i$。

目标函数仍为 2024--2030 年总利润最大化：
\[
\max Z=
\sum_{t\in T}\sum_{s\in S}\sum_{i\in C}P_{si}S_{tsi}
-
\sum_{t\in T}\sum_{l\in L}\sum_{k\in M_l}C_{lk}y_{tlk},
\]
其中，$C_{lk}$ 表示地块 $l$ 选择第 $k$ 个候选模式时的总种植成本。

在求解参数设置上，本文将最长求解时间设为 600 秒，相对最优间隙设为 2\%。
Gurobi 最终在 2.0431 秒内完成求解，模型状态为 OPTIMAL，目标函数值为 15503088.75 元，
MIPGap 为 1.4070\%，小于设定的 2\% 最优性间隙要求。因此，可以认为 Gurobi 在当前候选模式模型下得到了有效的优化结果。

\begin{table}[htbp]
\centering
\caption{Gurobi 对原始模型的计算结果}
\label{tab:gurobi-original}
\small
\setlength{\tabcolsep}{4pt}
\resizebox{\textwidth}{!}{%
\begin{tabular}{lrrrrrrc}
\toprule
算法 & 运行时间/s & 目标函数值/元 & 总收入/元 & 总成本/元 & 滞销浪费/斤 & MIPGap & 是否可行 \\
\midrule
Gurobi-原始模型 & 2.0431 & 15503088.75 & 19998908.75 & 4495820.00 & 3268800.00 & 1.4070\% & 是 \\
\bottomrule
\end{tabular}%
}
\end{table}

从表 \ref{tab:gurobi-original} 可以看出，Gurobi 得到的目标函数值为 15503088.75 元，
明显高于模拟退火算法和遗传算法所得结果。相较于遗传算法，Gurobi 的总利润更高，
同时滞销浪费产量由 7790400.00 斤下降至 3268800.00 斤，说明 Gurobi 在兼顾收益和市场销售上限方面表现更优。

进一步比较可知，模拟退火算法和遗传算法依赖随机扰动、交叉变异以及罚函数设计，
其结果受初始解和参数设置影响较大；而 Gurobi 通过数学规划方式在给定候选模式空间内进行系统搜索，
能够在较短时间内获得更高质量的解。因此，在本问题的原始模型求解中，Gurobi 结果可作为评价智能优化算法结果优劣的重要参照。

需要说明的是，Gurobi 输出结果中仍显示作物分散度存在一定超限情况。
这说明虽然模型已经满足不连续重茬和三年内豆类轮作等主要硬约束，但在田间管理便利性方面仍有进一步优化空间。
后续可以通过增加候选模式数量、调整作物分散度约束口径，或将作物分散度作为更严格的硬约束，
进一步提高种植方案的管理可操作性。

\section{原始模型的线性化改造与求解}

\subsection{线性化模型建立}

原始模型的非线性部分主要来自目标函数中的
$\min\{Q_{tsi},D_{si}\}$。为将模型转化为便于商业优化软件直接处理的混合整数线性规划模型，
引入实际销售量变量：
\[
S_{tsi}\geq 0,
\quad \forall t\in T,\ s\in S,\ i\in C.
\]

其中，$S_{tsi}$ 表示第 $t$ 年、第 $s$ 季作物 $i$ 的实际可销售产量。
在第（1）种情形下，实际销售量不能超过当季总产量，也不能超过预期销售量，因此增加如下线性约束：
\[
0\leq S_{tsi}\leq Q_{tsi},
\quad \forall t\in T,\ s\in S,\ i\in C,
\]
\[
0\leq S_{tsi}\leq D_{si},
\quad \forall t\in T,\ s\in S,\ i\in C.
\]

线性化后的目标函数为：
\[
\max Z_L=
\sum_{t\in T}\sum_{s\in S}\sum_{i\in C}P_{si}S_{tsi}
-
\sum_{t\in T}\sum_{l\in L}\sum_{s\in S}\sum_{i\in C}
K_{lsi}x_{tlsi}.
\]

同时，总产量约束仍为：
\[
Q_{tsi}=\sum_{l\in L}Y_{lsi}x_{tlsi},
\quad \forall t\in T,\ s\in S,\ i\in C.
\]

由于销售价格 $P_{si}\geq 0$，且模型目标为最大化总利润，在最优解中 $S_{tsi}$ 会尽可能取大。
在 $S_{tsi}\leq Q_{tsi}$ 和 $S_{tsi}\leq D_{si}$ 同时成立的条件下，最优时必有：
\[
S_{tsi}=\min\{Q_{tsi},D_{si}\}.
\]

因此，上述线性化处理不改变原始模型的经济含义和最优目标值，只是用线性约束显式刻画了
``实际销售量不超过产量和需求'' 这一关系。线性化后的完整模型由线性目标函数、总产量约束、
实际销售量约束以及前文给出的作物适宜性、面积、轮作、豆类作物和分散度约束共同组成，
属于混合整数线性规划模型。

\subsection{线性化模型的智能优化算法求解}

对于模拟退火算法和遗传算法而言，算法本身仍然直接搜索种植方案。
在评价一个候选方案时，先根据方案计算 $Q_{tsi}$，再由
\[
S_{tsi}=\min\{Q_{tsi},D_{si}\}
\]
计算实际销售量和销售收入。因此，线性化后的评价函数与原始模型中的评价函数在数值上完全一致。
本文采用与原始模型相同的算法参数和候选方案编码方式，对线性化模型进行重复评价。
计算结果如表 \ref{tab:linear-sa-ga-result} 所示。

\begin{table}[htbp]
\centering
\caption{线性化模型下两种智能优化算法的计算结果}
\label{tab:linear-sa-ga-result}
\small
\setlength{\tabcolsep}{3pt}
\resizebox{\textwidth}{!}{%
\begin{tabular}{lrrrrrrc}
\toprule
算法 & 运行时间/s & 原始利润/元 & 总收入/元 & 总成本/元 & 滞销浪费/斤 & 惩罚项 & 是否有效 \\
\midrule
模拟退火 SA & 9.1436 & 9705157.00 & 16368205.00 & 6663048.00 & 9168905.00 & 9000000.00 & 是 \\
遗传算法 GA & 8.6847 & 10921219.75 & 17233973.75 & 6312754.00 & 7790400.00 & 7000000.00 & 是 \\
\bottomrule
\end{tabular}%
}
\end{table}

从表 \ref{tab:linear-sa-ga-result} 可以看出，线性化模型下两种智能优化算法均可以得到有效结果。
由于线性化并未改变候选方案的利润计算方式，模拟退火和遗传算法所得利润、收入、成本及滞销浪费量
与原始模型评价结果一致。其中，遗传算法的利润为 10921219.75 元，高于模拟退火算法的
9705157.00 元，说明在相同候选空间和参数设置下，遗传算法仍具有更好的搜索效果。

\subsection{Gurobi 对线性化模型的求解}

在线性化模型中，Gurobi 不再需要使用一般约束表示最小值关系，而是直接处理变量
$S_{tsi}$ 及其上下界约束。这样得到的模型为标准混合整数线性规划模型，便于 Gurobi 调用
分支定界、割平面和启发式搜索等方法求解。

在候选模式编码下，线性化模型可写为：
\[
\max Z_L=
\sum_{t\in T}\sum_{s\in S}\sum_{i\in C}P_{si}S_{tsi}
-
\sum_{t\in T}\sum_{l\in L}\sum_{k\in M_l}C_{lk}y_{tlk},
\]
满足：
\[
\sum_{k\in M_l}y_{tlk}=1,
\quad \forall t\in T,\ l\in L,
\]
\[
Q_{tsi}=\sum_{l\in L}\sum_{k\in M_l}a_{lksi}Y_{lsi}A_l y_{tlk},
\quad \forall t\in T,\ s\in S,\ i\in C,
\]
\[
0\leq S_{tsi}\leq Q_{tsi},\quad
0\leq S_{tsi}\leq D_{si},
\quad \forall t\in T,\ s\in S,\ i\in C,
\]
以及原模型中的轮作、豆类作物和分散度等约束。

本文仍采用 600 秒时间限制和 2\% 相对最优间隙作为求解参数。
由于线性化模型与原始最小值模型等价，并采用相同候选模式空间，Gurobi 得到的目标函数值
与原始模型结果一致。计算结果如表 \ref{tab:gurobi-linear} 所示。

\begin{table}[htbp]
\centering
\caption{Gurobi 对线性化模型的计算结果}
\label{tab:gurobi-linear}
\small
\setlength{\tabcolsep}{4pt}
\resizebox{\textwidth}{!}{%
\begin{tabular}{lrrrrrrc}
\toprule
算法 & 运行时间/s & 目标函数值/元 & 总收入/元 & 总成本/元 & 滞销浪费/斤 & MIPGap & 是否可行 \\
\midrule
Gurobi-线性化模型 & 2.0431 & 15503088.75 & 19998908.75 & 4495820.00 & 3268800.00 & 1.4070\% & 是 \\
\bottomrule
\end{tabular}%
}
\end{table}

由表 \ref{tab:gurobi-linear} 可知，Gurobi 在线性化模型下能够得到有效解，
目标函数值为 15503088.75 元，MIPGap 为 1.4070\%，满足本文设定的 2\% 最优性间隙要求。
与智能优化算法相比，Gurobi 的目标函数值更高，滞销浪费产量更低，说明线性化后的数学规划模型
更适合使用商业优化软件进行系统搜索。

\subsection{线性化前后结果比较}

为比较原始模型和线性化模型的求解效果，将不同算法在两种模型表述下的主要结果汇总如表
\ref{tab:linear-compare} 所示。

\begin{table}[htbp]
\centering
\caption{原始模型与线性化模型计算结果比较}
\label{tab:linear-compare}
\small
\setlength{\tabcolsep}{3pt}
\resizebox{\textwidth}{!}{%
\begin{tabular}{llrrrrc}
\toprule
模型表述 & 算法 & 运行时间/s & 目标值或利润/元 & 总收入/元 & 滞销浪费/斤 & 是否有效 \\
\midrule
原始模型 & 模拟退火 SA & 9.1436 & 9705157.00 & 16368205.00 & 9168905.00 & 是 \\
线性化模型 & 模拟退火 SA & 9.1436 & 9705157.00 & 16368205.00 & 9168905.00 & 是 \\
原始模型 & 遗传算法 GA & 8.6847 & 10921219.75 & 17233973.75 & 7790400.00 & 是 \\
线性化模型 & 遗传算法 GA & 8.6847 & 10921219.75 & 17233973.75 & 7790400.00 & 是 \\
原始模型 & Gurobi & 2.0431 & 15503088.75 & 19998908.75 & 3268800.00 & 是 \\
线性化模型 & Gurobi & 2.0431 & 15503088.75 & 19998908.75 & 3268800.00 & 是 \\
\bottomrule
\end{tabular}%
}
\end{table}

从表 \ref{tab:linear-compare} 可以看出，原始模型和线性化模型在目标值上保持一致。
其原因在于线性化模型通过引入 $S_{tsi}$ 并施加
$S_{tsi}\leq Q_{tsi}$、$S_{tsi}\leq D_{si}$，等价地表达了原始模型中的
$\min\{Q_{tsi},D_{si}\}$。因此，对于同一套种植方案，两种模型表述计算出的销售收入和利润相同。

不过，线性化改造仍然具有重要意义。对于智能优化算法，线性化后的模型使适应度函数的含义更加清楚，
便于将实际销售量、总产量和预期销售量分别统计；对于 Gurobi，线性化模型将原来的最小值关系转化为
标准线性约束，使模型成为混合整数线性规划模型，表达更加规范，也更便于商业优化软件稳定求解。

\section{结论}

本文针对 2024 年全国大学生数学建模竞赛 C 题问题一第（1）种情形，先建立了包含
$\min\{Q_{tsi},D_{si}\}$ 的原始优化模型，并分别采用模拟退火算法、遗传算法和 Gurobi 进行求解。
随后，引入实际销售量变量 $S_{tsi}$，通过线性约束
$S_{tsi}\leq Q_{tsi}$ 和 $S_{tsi}\leq D_{si}$ 对原始模型进行线性化改造。

计算结果表明，模拟退火算法、遗传算法和 Gurobi 在线性化模型下均可以得到有效结果。
由于线性化模型与原始模型在数学上等价，智能优化算法在线性化前后的计算结果基本一致。
在三种算法中，Gurobi 得到的目标函数值最高，为 15503088.75 元，运行时间为 2.0431 秒，
同时滞销浪费产量最低。因此，在本文设定的候选模式空间和参数条件下，Gurobi 线性化模型结果
可作为问题一第（1）种情形的最优有效方案。

\end{document}
